\section*{Chapter \ref{ch:g12:diffeq} --- Differential Equations}

\begin{solution}{exo:g12:diffeq:1}
\emph{(a)} $y = C\eu^{3x}$; $y(0) = 2$ gives $y = 2\eu^{3x}$.

\emph{(b)} $y' = -\frac12 y$, so $y = C\eu^{-x/2}$; $y(0) = -1$ gives
$y = -\eu^{-x/2}$.

\emph{(c)} Equilibrium $y_p = 5$; $y = C\eu^{-x} + 5$; $y(0) = 0$ gives
$C = -5$: $y = 5\left(1 - \eu^{-x}\right)$.
\end{solution}

\begin{solution}{exo:g12:diffeq:2}
$N' = kN$, so $N(t) = N_0 \eu^{kt}$. Doubling in $3$ hours means
$\eu^{3k} = 2$, \emph{i.e.}
\[
k = \frac{\ln 2}{3} \approx 0.231\ \text{h}^{-1}.
\]
\end{solution}

\begin{solution}{exo:g12:diffeq:3}
$y_p'(x) = \eu^x + x\eu^x = y_p(x) + \eu^x$: $y_p$ is a particular solution
(this is the resonant case of \cref{met:g12:diffeq:particular}). By
\cref{prop:g12:diffeq:general}, the solutions are
$y = C\eu^{x} + x\eu^{x} = (C + x)\,\eu^x$, $C \in \R$.
\end{solution}

\begin{solution}{exo:g12:diffeq:4}
Try $y_p = \beta\eu^x$: $\beta\eu^x = -2\beta\eu^x + \eu^x$ gives
$3\beta = 1$, so $y_p = \frac13\eu^x$. General solution
$y = C\eu^{-2x} + \frac13\eu^{x}$; the condition $y(0) = 1$ gives
$C = \frac23$:
\[
y(x) = \frac{2}{3}\,\eu^{-2x} + \frac{1}{3}\,\eu^{x}.
\]
\end{solution}

\begin{solution}{exo:g12:diffeq:5}
$N(t) = N_0\,\eu^{-\lambda t}$ with $\lambda = \frac{\ln 2}{5730}$
(\cref{ex:g12:exp:halflife}). We solve $\eu^{-\lambda t} = 0.6$:
\[
t = \frac{\ln(1/0.6)}{\lambda} = 5730\,\frac{\ln(5/3)}{\ln 2}
\approx 5730 \times \frac{0.5108}{0.6931} \approx 4220 \text{ years}.
\]
\end{solution}

\begin{solution}{exo:g12:diffeq:6}
\emph{1.} Salt flows in at $0.2 \times 5 = 1$ kg/min. The outflow carries
concentration $\frac{m}{100}$ kg/L at $5$ L/min, \emph{i.e.} $\frac{m}{20}$
kg/min. Hence $m' = 1 - \frac{m}{20}$.

\emph{2.} Equilibrium $m_p = 20$; $m(t) = 20 + C\eu^{-t/20}$, and
$m(0) = 0$ gives $C = -20$:
\[
m(t) = 20\left(1 - \eu^{-t/20}\right) \xrightarrow[t\to+\infty]{} 20 \text{ kg}.
\]
In the long run the tank's concentration equals that of the incoming brine:
$0.2$ kg/L $\times$ $100$ L $= 20$ kg.
\end{solution}

\begin{solution}{exo:g12:diffeq:7}
\emph{1.} $v' = g - \frac{k}{m} v$ with $\frac km = \frac{16}{80} = 0.2$.
Equilibrium $v_\infty = \frac{mg}{k} = \frac{9.8}{0.2} = 49$ m/s;
$v(t) = 49 + C\eu^{-0.2t}$, and $v(0) = 0$ gives
\[
v(t) = 49\left(1 - \eu^{-0.2 t}\right).
\]

\emph{2.} Terminal velocity $49$ m/s ($\approx 176$ km/h). We want
$1 - \eu^{-0.2t} = 0.95$, \emph{i.e.} $\eu^{-0.2t} = 0.05$:
\[
t = \frac{\ln 20}{0.2} \approx \frac{3.00}{0.2} \approx 15 \text{ s}.
\]
\end{solution}

\begin{solution}{exo:g12:diffeq:8}
\emph{1.} $z = \frac1y$ gives $z' = -\frac{y'}{y^2}
= -\frac{y(1-y)}{y^2} = -\frac{1-y}{y} = -\frac1y + 1 = -z + 1$.

\emph{2.} Equilibrium $z_p = 1$, so $z(t) = 1 + C\eu^{-t}$. From
$y(0) = \frac{1}{10}$, $z(0) = 10$, so $C = 9$ and
\[
y(t) = \frac{1}{1 + 9\,\eu^{-t}} .
\]

\emph{3.} As $t \to +\infty$, $9\eu^{-t} \to 0$ and $y(t) \to 1$. The curve
is the classic S-shaped (\emph{sigmoid}) logistic curve: slow growth at
first ($y$ small, $y' \approx y$), fastest growth when $y = \frac12$
(where $y' = y(1-y)$ is maximal), then saturation towards the carrying
capacity $1$.
\end{solution}

\begin{solution}{exo:g12:diffeq:9}
By \cref{thm:g12:diffeq:affine}, $y(x) = C\eu^{ax} - \frac ba$. Hence
\[
u_{n+1} = C\eu^{a(n+1)} - \frac ba
= \eu^{a}\left(C\eu^{an} - \frac ba\right) + \frac ba\left(\eu^a - 1\right)
= q\,u_n + r
\]
with $q = \eu^{a}$ and $r = \frac{b}{a}\left(\eu^{a} - 1\right)$.
The condition $\abs q < 1$ means $\eu^a < 1$, \emph{i.e.} $a < 0$: exactly
the condition under which the solutions of the \omterm{def:g12:diffeq:ode}{differential equation}
converge to the equilibrium $-\frac ba$ --- and indeed the \omterm{pb:g10:functions:1}{fixed point} of
the recurrence is $\frac{r}{1 - q} = -\frac{b}{a}$.
\end{solution}

\begin{solution}{pb:g12:diffeq:1}
\textbf{1.} $y = 2\eu^{3t}$. For the second: equilibrium $3$,
so $y = 3 + C\eu^{-2t}$ with $y(0) = 0$: $C = -3$:
$y = 3\left(1 - \eu^{-2t}\right)$.

\textbf{2.} $\left(y\eu^{-at}\right)' = y'\eu^{-at} -
ay\eu^{-at} = (y' - ay)\eu^{-at} = 0$: the product is a
constant $C$, so $y = C\eu^{at}$ --- no solution escapes.

\textbf{3.} Equilibrium: $y \equiv 3$. Every other solution is
$3 + C\eu^{-2t}$ with $C \neq 0$: the exponential dies and the
solution glides to $3$ --- a stable \omterm{pb:g10:functions:1}{fixed point}, the
\omterm{def:g12:limcont:continuity}{continuous} twin of the recurrence \omterm{pb:g10:functions:1}{fixed points} of
\cref{pb:g11:seq:1}.

\textbf{4.} $y = y_0\eu^{-t/\tau}$ halves when
$\eu^{-t/\tau} = \frac12$: $t = \tau\ln 2$.

\textbf{5.} All solutions are vertical translates of the decay
towards $3$: those starting above fall to $3$, those below
rise to $3$, and the constant solution $3$ sits still --- a
funnel around the equilibrium.

\textbf{6.} $\lambda = \frac{\ln 2}{5730} \approx 1.21 \times
10^{-4}$ per year.

\textbf{7.} $\eu^{-\lambda t} = 0.2$:
$t = \frac{\ln 5}{\lambda} \approx 13\,300$ years.

\textbf{8.} $t = \frac{\ln(1/0.15)}{\lambda} \approx 15\,700$
years: the bulls of Lascaux are late Ice Age.

\textbf{9.} $t = \frac{\ln(1/0.53)}{\lambda} \approx 5\,250$
years: within a lifetime of the archaeologists' value ---
the clock works.

\textbf{10.} Ten half-lives leave $2^{-10} \approx 0.1\,\%$:
at the edge of measurement. After $65$ million years the
fraction is $2^{-65\,000\,000/5\,730} \approx 2^{-11\,344}$:
no atom of the original stock remains in any fossil ---
dinosaurs are dated by slower clocks (potassium--argon and
kin).

\textbf{11.} $52\,\%$ gives $\approx 5\,405$ years,
$54\,\%$ gives $\approx 5\,095$: the $\pm 1\,\%$ of chemistry
becomes roughly $\pm 155$ years of history --- precision in
the lab is precision in the museum label.

\textbf{12.} $z = T - T_a$ satisfies $z' = -kz$:
$z = (T_0 - T_a)\eu^{-kt}$, and $T = T_a + z$.

\textbf{13.} The initial gap is $T_0 - T_a = 70$; after five
minutes the gap is $70 - 20 = 50$, so $50 = 70\eu^{-5k}$:
$k = \frac{\ln(70/50)}{5} \approx 0.067$ per minute. Drinkable:
$55 - 20 = 35 = 70\eu^{-kt}$:
$t = \frac{\ln 2}{k} \approx 10.3$ minutes.

\textbf{14.} The loss rate is proportional to the gap
$T - T_a$: piping-hot coffee bleeds heat fastest. Adding the
milk \emph{now} slashes the gap immediately, so less heat is
lost over the ten minutes; adding it at the end lets the
coffee cool at full speed first. For the warmest cup: milk
first. (Impatient hosts have the thermodynamics backwards.)

\textbf{15.} Gaps from ambient: at midnight $10$, an hour
later $8$: $\eu^{-k} = 0.8$:
$k = \ln\frac{10}{8} \approx 0.223$ per hour. At death the gap
was $17$; it decayed to $10$ by midnight:
$s = \frac{\ln(17/10)}{k} \approx 2.4$ hours. Time of death:
about $21{:}40$.

\textbf{16.} A heated or draughty room ($T_a$ not constant)
bends the clock either way; clothing or body mass changes $k$
(the coroner's tables adjust for it) --- a wrong $k$ scales
the whole \omterm{def:g10:numbers:interval}{interval}; and a body moved from elsewhere resets
$T_a$ mid-decay, faking an earlier or later death. The
\omterm{def:g10:algebra:equation}{equation} is honest; the assumptions carry the risk.

\textbf{17.} $z = \frac1y$ obeys $z' = 1 - z$:
$z = 1 + 9\eu^{-t}$ (from $z(0) = 10$), so
$y = \frac{1}{1 + 9\eu^{-t}}$. Inflection at $y = \frac12$:
$9\eu^{-t} = 1$: $t = \ln 9 \approx 2.2$ --- the epidemic's
turning date, straight from the formula.

\textbf{18.} Terminal velocity: $v' = 0$ at $v = 20$ m/s.
Solution: $v(t) = 20\left(1 - \eu^{-t/2}\right)$. Then
$0.95$: $\eu^{-t/2} = 0.05$: $t = 2\ln 20 \approx 6$ s ---
raindrops hit their cruising speed within seconds, which is
why rain does not kill.

\textbf{19.} Steady state: $c = 8$. Half of it:
$4 = 8\left(1 - \eu^{-t/2}\right)$: $\eu^{-t/2} = \frac12$:
$t = 2\ln 2 \approx 1.4$ hours. The drip is the \omterm{def:g12:limcont:continuity}{continuous}
twin of the loan: constant inflow, proportional outflow ---
\cref{exo:g12:diffeq:9} makes the dictionary exact.

\textbf{20.} Decay ($a < 0$, $b = 0$), cooling ($y = T - T_a$),
falling ($a < 0$, $b > 0$: approach to equilibrium from
below), dosing (same, in a vein): one \omterm{def:g10:algebra:equation}{equation}, four worlds.
The loop: state the law of change; write the \omterm{def:g10:algebra:equation}{equation}; solve
(\cref{met:g12:diffeq:solve}); calibrate constants on
measurements; predict; then audit the assumptions (question
16). Oscillation needs $y'' = -\omega^2 y$ --- met, solved and
set swinging in \cref{pb:g12:trigo:1}.
\end{solution}
